\chapter{The experimental protocol and the provenance pipeline}
\label{ch:protocol}

A number in the paper is a macro. The macro was written by a generator script that read
a CSV file. The CSV file was written by an experiment script under a frozen
configuration, on a machine whose identity and code revision it recorded, and it was
then pinned by a cryptographic hash in a registry entry that had been written before the
experiment ran. Two independent checks stand between the CSV and the PDF: the manuscript
check regenerates every generated fragment and refuses any typed decimal, and an audit
script recomputes every headline number from the CSVs with independent code. This
chapter walks that chain link by link, defines every object in it, and explains the
rules the project follows when something has to change. It is the chapter that lets a
reader who is not the author decide how much to trust a number.

\section{Two frozen objects}
\label{sec:frozen}

The protocol separates two things that earlier versions of this work had mixed: the
\emph{world} the experiments are run in, and the \emph{comparison} the experiments make.
Each is one Python dataclass whose fields cannot be mutated after construction; an
experiment that needs a different value must say so at the call site, and the
difference is then recorded in its provenance.

\subsection{The world: \code{FrozenConfig}}

\code{experiments/frozen\_config.py} defines \code{FROZEN}, an instance of
\code{FrozenConfig}. Field by field:
\begin{description}[leftmargin=2em,itemsep=0.2em]
\item[\code{rho}] the correlation of the chain, $\corr=0.85$; one value everywhere.
\item[\code{n\_sites}] the chain length $\nsites=32$, which is also the ambient dimension.
\item[\code{innovation}] the default innovation family, Laplace (excess kurtosis $3$; its
  third cumulant vanishes, so the fourth is the leading non-Gaussian coordinate).
\item[grid] the fields \code{n\_grid} and \code{half\_width}: the quadrature grid, $\ngrid=401$ points on
  $[-\halfwidth,\halfwidth]=[-8,8]$ (\Cref{ch:grid-quadrature}); the field's docstring
  records the boundary-mass statistics that justified it and a warning about quoting a
  quantile as if it were a bound.
\item[\code{t\_grid}] the twelve-level schedule, $\log$-spaced on $[0.05,3]$ and rounded to
  four decimals (\Cref{ex:schedule} below).
\item[\code{n\_seeds}] sixteen replicate seeds, with the reason (three seeds once read a
  falling curve as flat).
\item[archival sizes] the fields \code{sizes} and \code{efficiency\_sizes}: the archival size lists of the earlier
  comparison; the registered experiments take their sizes from the protocol instead.
\item[archival held-out] the fields \code{n\_heldout} and \code{heldout\_seed\_offset}: the archival held-out design.
\item[EM stopping] the fields \code{em\_max\_iters}, \code{em\_loglik\_tol} and \code{em\_shape\_tol}: EM's archival
  cap, its likelihood tolerance of $10^{-9}$ nats per transition (the one the registered
  experiments still pass to \code{fit\_em}), and the shape tolerance, with a docstring
  explaining why convergence is judged on the shape and never on $\corr$ alone.
\item[\code{n\_components}] the default mixture size $\ncomp=8$.
\item[\code{ring\_*}] the fields of an exploratory model (a ring of frames) that the
  paper does not use.
\item[\code{base\_seed}] the seed from which \code{seeds()} derives the replicate seeds.
\end{description}
Three methods matter: \code{innovation\_variance} returns $\ivar=1-\corr^2$, the
variance-matching choice that makes every family share $\Sig_0$; \code{narrowed(**kw)}
returns a copy with some fields changed, for a deliberate and reported deviation;
\code{as\_dict()} is what every provenance record stores. The module also defines
\code{config\_hash}, a twelve-character prefix of the SHA-256 digest of the sorted
\code{key=value} lines of the configuration, and \code{provenance(config)}, which records
the configuration, its hash, the frozen hash, a flag \code{is\_frozen} and the dictionary
of deviations from \code{FROZEN}.

\begin{example}[The twelve-level schedule, by hand]
\label{ex:schedule}
\code{\_log\_grid(0.05, 3.0, 12)} returns
$t_k=0.05\cdot60^{(k-1)/11}$ for $k=1,\dots,12$, since $3/0.05=60$. The ratio between
consecutive levels is $60^{1/11}=e^{(\log60)/11}=e^{4.0943/11}=e^{0.3722}\approx1.4509$.
Multiplying out: $0.0500$, $0.0725$, $0.1053$, $0.1527$, $0.2216$, $0.3215$, $0.4665$,
$0.6768$, $0.9819$, $1.4247$, $2.0671$, $3.0000$. These are the levels every table and
figure of the paper is pooled over, and the reason the network's time features take
$\log t$ as their argument (\Cref{ch:neural-baselines}). The crossover
$\tstar=\tfrac12\log\lmax(\Sig_0)$ at $\corr=0.85$ falls between the ninth and tenth level.
\end{example}

\subsection{The comparison: \code{Protocol}}

\code{experiments/protocol.py} defines \code{PROTOCOL}, version 2, the registered
protocol of the cluster runs. Field by field:
\begin{description}[leftmargin=2em,itemsep=0.2em]
\item[\code{sizes}] the training-set sizes $\nseq\in\{32,128,512,2048\}$.
\item[seeds] the fields \code{seeds} and \code{dev\_seeds}: the confirmatory seeds $0$--$15$ and the
  development seeds $10000$--$10003$ (\Cref{ch:model-selection}).
\item[bundle sizes] the fields \code{n\_val} and \code{n\_test}: $1024$ validation and $2048$ test chains.
\item[\code{bundle\_tags}] the stream tags \code{exp31-train}, \code{exp31-val},
  \code{exp31-test}; the training stream is keyed by \code{(tag, seed, M)} and the other
  two by \code{(tag, seed)}, so that every experiment at the same key sees identical
  chains (\Cref{sec:streams}).
\item[\code{estimand}] the schedule-pooled relative score error, written out as a
  formula in the field's value so that it is part of the hash.
\item[regions] the fields \code{headline\_region}, \code{selection\_region} and
  \code{bulk\_margin}: all sites for the headline, the bulk for selection, four sites
  trimmed at each end.
\item[\code{capacities}] $\ncomp\in\{1,2,4,8,16\}$.
\item[EM arm] the fields \code{em\_checkpoints}, \code{em\_init\_rho},
  \code{em\_init\_var} and \code{one\_view}: the EM ladder to $1200$, the initialisation
  $\corr=0.3$ and innovation variance $0.8$, and the one-view flag.
\item[network arms] the fields \code{net\_steps}, \code{net\_checkpoints},
  \code{batch\_size}, \code{lrs}, \code{parameterizations} and \code{grad\_clip}: the
  network recipe of \Cref{ch:neural-baselines}.
\item[screen] the fields \code{screen\_sizes}, \code{screen\_restarts} and
  \code{screen\_steps}: the screen budget as registered; E03's amendment narrowed it
  before the screen ran.
\item[\code{factors}] the table of which experiment varies which single factor, with
  everything else held at the protocol: the \term{ceteris paribus} design.
\end{description}
\code{V1} is the laptop protocol of the first runs (ladder to $400$, capacities to $8$,
two development seeds, one restart of $16{,}000$ steps), kept so that a version-1 output
can be described exactly and so that \code{v1\_overrides()} can reproduce it;
\code{protocol\_hash()} digests the protocol together with the world. A params file that
records the protocol version and the resolved settings lets a reader tell which protocol
produced a file without remembering anything.

\subsection{Keyed random streams}
\label{sec:streams}

\begingroup\sloppy
Every random draw in the registered experiments comes from
\code{rng\_for(*keys)} (\code{src/markov\_score/utils.py}): the keys, prefixed by a
project constant, are hashed by a stable function into a NumPy \code{SeedSequence}, from
which a generator is created. The same keys give the same stream in every process and on
every machine. This is the \term{common random numbers} device of simulation: two arms
that call \code{rng\_for("exp31-train", seed, M)} see the same training chains, and a
difference between them cannot be Monte Carlo noise. The chapter on model selection
explains how the tags separate training, validation and test data; the negative-results
record explains how a driver that re-created its generator inside a loop defeated the
device (\Cref{sec:amendments}).\par\endgroup

\section{The registry}
\label{sec:registry}

\subsection{One JSON entry per experiment}

Every registered experiment has a file \code{experiments/registry/<id>.json}. Its fields:
\begin{description}[leftmargin=2em,itemsep=0.2em]
\item[identity] the fields \code{id}, \code{script} and \code{script\_sha256}: the experiment's identifier, its
  driver under \code{experiments/}, and the driver's SHA-256 at registration, so a later
  edit of the script is detectable.
\item[\code{hypothesis}] the statement the experiment tests, written before it runs, with
  its prediction.
\item[\code{design}] (or \code{parts} and \code{settings}) the factor varied, what is
  held, the sizes, the arms, the estimator, the grid.
\item[\code{predeclared\_analysis}] the generator that will read the output and the
  statistic it will compute, named before the data exist, including what will be
  reported as censored or bracketed rather than fitted.
\item[\code{seeds}] the development and confirmatory seed blocks.
\item[\code{status}] a string from the vocabulary \textsc{pre-registered},
  \textsc{exploratory}, \textsc{confirmatory}, \textsc{done}, \textsc{withdrawn}, with a
  \textsc{frozen} prefix added by the freezing tool and the cluster job identifiers
  appended as the run progresses.
\item[history] the fields \code{registered} and \code{amendments}: the registration date and a list of dated
  amendments, each with what changed and why (\Cref{sec:amendments}).
\item[\code{hpc\_runs}] the cluster jobs, their deployments and their fates.
\item[free text] the fields \code{motivation}, \code{limitations} and \code{outcome}, present where the
  experiment was registered during the audit or where it returned a result against its
  hypothesis.
\item[freeze record] the fields \code{frozen\_outputs}, \code{frozen\_at},
  \code{frozen\_commit} and \code{frozen\_run\_commits}, written by the freezing tool
  (\Cref{sec:freeze}).
\end{description}
\code{docs/EXPERIMENT\_REGISTRY.md} is the human-readable table of the same entries;
\code{tools/registry.py --sync} writes each output's commit back into it, and
\code{--check} fails if any experiment with rows lacks a clean provenance record or if a
frozen file changed.

\subsection{Pre-registration}

\term{Pre-registration} means that the hypothesis, the design, the seeds and the analysis
are written down and committed before the confirmatory data exist. Its purpose is to
make the analysis independent of the outcome: a choice made after seeing the data (which
levels to fit a slope on, which arm to call ``matched'', which seeds to report) can be
tuned to the result, and a choice made before cannot. The paper's contract states the
rule in four lines: confirmatory settings live in the protocol and the registry entry,
frozen before the first confirmatory cell runs; development seeds never enter a
reported number; confirmatory seeds never enter a selection; a running confirmatory
experiment is never changed.

\subsection{Pre-data amendments}
\label{sec:amendments}

A registration can be wrong before it runs, and the instrument for correcting it is the
\term{pre-data amendment}: a dated entry in the registry's \code{amendments} list, made
before any confirmatory cell has produced a number, stating what changed and why. The
registry holds several, and they are worth reading as specimens of the rule:
\begin{itemize}
\item E03's screen budget was cut from four development seeds, three restarts and
  $32{,}000$ steps to two seeds, one run and $16{,}000$ steps after the single-thread cost
  was measured, with the confirmatory protocol untouched;
\item E04's held-out bundle was reduced from $2048$ to $512$ chains after the finest grid
  produced no completed cell in nineteen hours, with the estimand (a paired movement
  between resolutions) unchanged;
\item E08's site-varying strength $0.2$ was replaced by $0.14$ when the construction was
  found invalid at $\corr=0.85$ (a site coefficient exceeds one), with the invalid cells
  recorded as never having existed;
\item E12b's analysis plan was changed from ``compare the fitted slopes with $4$ and $8$''
  to ``compare $M^\ast(t)$ with the exact information curve'' after E21 showed that the
  asymptotic exponents are not the prediction on the swept levels, while the job was
  still pending in the queue;
\item E14 was amended twice pre-data: to record the whole held-out implicit objective
  curve so that a diverged fit can be called a divergence, and, on 20 September, to fix
  and guard its data generator.
\end{itemize}

The last case is the most instructive and is recorded as entry N-22 of
\code{docs/NEGATIVE\_RESULTS.md}. E14's first screen gave the window head, at the frozen
point and with E03's own configuration, a test risk several times the value the same
network had reached in E03's matched arm, with every held-out objective curve rising
from the first checkpoint. The first diagnosis blamed the divergence estimator (the
screen had used Hutchinson's one probe where E03 used the exact trace); the driver was
changed to E03's per-arm estimators and learning rates, and a re-run failed the same
way. What settled it was running E03's training function and E14's on the same data:
both gave the bad number on E14's data and the good number on E03's, so the cause was
the data. In the driver's \code{\_common} the training chains had been drawn by a
comprehension that re-created the generator \emph{inside} the loop,
\begin{center}
\code{prior.sample(rng\_for("e14-train", seed, L, rho), L) for \_ in range(n)},
\end{center}
so all $2048$ training chains were one and the same draw, and both arms had fitted two thousand
copies of a single chain. The fix creates one generator per set and asserts that the
chains are distinct (\code{assert len(np.unique(A[:, 0])) > 1}); with it E14's pipeline
reproduces E03's value at the frozen point. The void screens and confirmations are kept
under \code{outputs/E14\_screen\_invalid/}, and a search of the current drivers found no
other instance of the pattern.

Why is a pre-data amendment the right instrument for such a fix, rather than a silent
edit or a fresh registration? Four reasons, in the order the contract gives them.
First, the guarantee that pre-registration buys is that the confirmatory design and
analysis were not chosen in the light of a confirmatory result. A runs whose data were
degenerate produced no confirmatory result about the registered question at all, so
amending before any valid cell exists preserves the guarantee exactly; the amendment's
timestamp, placed against the job log, is the evidence. Second, the alternative of
editing the driver silently would leave the registry describing a run that never
happened as registered, and the alternative of a new registration would suggest that the
scientific design had changed, when only its implementation had; the amendment says
precisely which. Third, the rule that a frozen or running output is never overwritten
still applies: the void runs are named, kept and excluded, so a reader can see what was
measured, why it is void, and that nothing was quietly replaced. Fourth, the amendment
converts the lesson into code, an assertion that a training set has as many distinct
chains as it claims, in the same spirit as the typed-literal gate and the
dirty-tree refusal: a defect that was caught once by a person is caught next time by the
machine. The audit's own summary of the lesson is the sentence to remember: a fitted
parameter that contradicts the generating law is a data check, not an estimator result.

\section{What a run records: the provenance record}
\label{sec:provenance}

Every part of every experiment writes \code{params\_<part>.json} beside its CSVs
(\code{provenance()} in \code{experiments/common.py}). Its fields:
\begin{description}[leftmargin=2em,itemsep=0.2em]
\item[machine] the fields \code{timestamp}, \code{hostname}, \code{platform},
  \code{python}, \code{numpy}, \code{cpu\_count}, \code{blas\_threads},
  \code{slurm\_job\_id} and \code{slurm\_array\_task\_id}: where and when the run happened
  and with what.
\item[revision] the fields \code{git\_commit}, \code{git\_branch} and \code{git\_dirty}:
  the revision, read from \code{git} when a repository is present; \code{git\_dirty} is
  the porcelain status, so a non-empty value means the commit alone does not reproduce
  the run.
\item[archive] the fields \code{source\_archive\_sha256}, \code{revision\_source} and
  \code{source\_is\_clean}:
  on a cluster node without \code{.git}, the commit and the digest of the exact source
  archive that was shipped are read from a \code{REVISION} file written at deployment
  (\Cref{sec:hpc}); \code{revision\_source} says which of \code{git}, the file or the
  environment supplied them. The archive digest exists because a stamp and a separately
  synchronised tree once disagreed and an output named a commit at which its own command
  could not have run.
\item[configuration] the field \code{resolved\_config\_hash}, a digest of the settings the run actually used,
  overrides included, so two runs of the same commit under different \code{--set} values
  are distinguishable.
\item[settings] the fields \code{settings}, \code{quick}, \code{workers} and
  \code{frozen}: the resolved settings dictionary, whether the run was a smoke test, the
  worker count, and the world.
\end{description}
With \code{NGBP\_REQUIRE\_CLEAN=1} in the environment a run refuses to start unless it can
name a commit, carries an archive digest and reports a clean tree. \term{Dirty} is
judged on the \emph{code path} only: \code{tools/provenance\_gate.py} computes the static
import closure of the driver and counts an uncommitted file as contaminating only if it
lies in that closure, so an edit to an unrelated script does not void a run and an edit
to a module the run imported does.

\section{Resumable cells}
\label{sec:cells}

Every multi-cell driver runs through \code{run\_cells} in \code{experiments/common.py}.
A \term{cell} is one unit of work (one seed at one size, say), identified by a key made
from its arguments. The worker that computes a cell writes that cell's rows to
\code{<output-dir>/cells/<csv stem>/<key>.json} by an atomic rename, and the driver
rebuilds the CSV from the cell files after every completion. Cells whose file exists are
skipped, so an interrupted run resumes by reissuing the same command; a cell that raised
leaves \code{<key>.failed} with its traceback and stops the run; duplicate keys are
refused. Nothing but a short key travels through the process pool, because a pool that
returned multi-kilobyte results once stopped delivering them after fifteen cells while
its workers computed on for five hours; the driver detects completion from the files.
The design is what allows a cluster job that hit its time limit to be resubmitted into
the same output directory.

\section{Freezing}
\label{sec:freeze}

\begin{definition}[Frozen output]
\label{def:frozen}
A registered output is \term{frozen} when every file of it that the manuscript may read
(\code{*.csv}, \code{params\_*.json}, \code{winners.json}, \code{floor\_fits.json}) has
been hashed with SHA-256 into its registry entry's \code{frozen\_outputs}, together with
the time, the repository commit at freezing and the commits recorded in the run's params
files.
\end{definition}

\code{tools/freeze.py <id>} does this and refuses a directory whose params record a
dirty code path or no commit. \code{tools/freeze.py --check} recomputes every digest and
fails on any byte change; it runs inside the registry check and inside the manuscript
check. \term{SHA-256} is a cryptographic hash: a fixed function from a file's bytes to
a $256$-bit digest (sixty-four hexadecimal characters) such that no feasible computation
produces two files with the same digest; a single changed byte changes the digest
unpredictably. Freezing therefore turns ``the file the paper cites'' into a value that
can be checked rather than a name that can be reassigned. The rules that follow from it:
a frozen file is immutable; a re-run is a new artefact under a new registry key or a
versioned directory, never over the old one; and a generator's source must be frozen
before its numbers may enter the manuscript.

\section{From a CSV to a macro: the generators}
\label{sec:generators}

\begingroup\sloppy
Every number, table and figure in the paper is produced by a script under
\code{tools/} that reads frozen outputs and writes either a fragment
\code{paper/sections/<name>.tex} or a figure \code{paper/figures/<name>.pdf}. A number
fragment is a list of \code{\textbackslash newcommand} definitions, one per quantity,
with a header line saying it was generated. The macro names are built by a fixed
encoding: a prefix for the experiment (\code{cap}, \code{base}, \code{learn},
\code{mis}, \code{aud}, \ldots), a quantity, and the digits of any level or strength
spelt out with \code{.} written as \code{p}, so that the name is a legal LaTeX command;
sizes are lettered $32\to$ A, $128\to$ B, $512\to$ C, $2048\to$ D. For instance the
root-mean-square error of $\hat\corr$ at level $t=0.4$ from E06 is the macro
\code{\textbackslash learnRmseRhoTzeropfour}, and the level itself is emitted as
\code{\textbackslash learnLevelTzeropfour}, so that the prose never types a level next to
a number computed at it. \code{paper/main.tex} inputs every fragment, and the body files
use only the macros.\par\endgroup

Generators refuse rather than warn: on a missing input, on a params file with a dirty
code path, on an incomplete design (a confirmatory seed missing a capacity), on an
unfrozen source, and in one case on a wrong constant (the theory generator refuses to
emit if the measured closure ratio for the Laplace chain lies outside a declared window
around one). A cell whose narrowest fitted mixture component is under two grid cells is
flagged with a dagger and reported, not dropped; the paper's protocol appendix now says
so, after the audit found an earlier sentence that said ``refused''.

\section{The manuscript check and the independent audit}
\label{sec:checks}

\subsection{\code{tools/check\_manuscript.py}}

The manuscript check runs five gates and exits non-zero on any failure:
\begin{enumerate}
\item \emph{Regenerate and diff.} Every generator (a fixed list plus an optional list of
  generators keyed by the CSV they need, skipped if that CSV is absent) is run into a
  temporary tree that links to everything except \code{paper/sections}, and each emitted
  fragment is compared line by line with the committed one. A difference means an input
  or a generator changed without the fragment being rebuilt, or that a fragment was
  hand-edited; this gate exists because a generated table once shipped with numbers from
  a superseded run beside a paragraph that carried the frozen ones.
\item \emph{Typed literals.} Every decimal number, every number in exponent notation and
  every $\times10^{k}$ in \code{main.tex}, the body files and the non-generated sections
  is matched by a regular expression, skipping lines that define commands, labels,
  references, citations or inputs; a match must appear in \code{paper/LITERALS\_ALLOWED.txt}
  with a justification (an analytic constant, a definition of the protocol, a quoted
  prior-literature value), or the check fails and names the line. This is what makes
  ``every number is a macro'' a property of the build rather than a promise.
\item \emph{Unresolved references.} The build log is scanned for undefined references
  and citations and for multiply defined labels.
\item \emph{Frozen sources.} \code{freeze.py --check} runs, and any CSV that exists,
  feeds a generator and is not frozen fails the check.
\item \emph{Overfull boxes.} The paper is built with \code{tectonic}; any overfull
  horizontal box in a generated table or a figure file fails, and any box more than ten
  points too wide anywhere fails (\code{docs/FIGURE\_TABLE\_STANDARDS.md}, rules T1 and F4).
\end{enumerate}

\subsection{\code{tools/audit\_claims.py}}

The generators are the only path by which a number reaches the manuscript, which is a
strength for provenance and a weakness for auditing: a bug in a generator is invisible to
the manuscript check, which only asks whether the macros match what the generator emits.
The audit script recomputes the headline quantities a second way, in short explicit
blocks of plain \code{pandas} on the frozen CSVs with no shared helpers, reads the
macros directly from \code{paper/sections/*-numbers.tex}, and compares. A comparison
passes when the recomputed value is within a tolerance of $0.0015$, or within half a unit
of the macro's last printed decimal, whichever is larger; a block that cannot run is
itself a finding. Blocks exist for the capacity table, the baselines and the matched
arm (with its own percentile bootstrap), the length study, generation, the
misspecification brackets, the regime sweep, the grid envelope, the learning
experiment's at-cap rates, the phase ladder (including a re-pooling of the per-level rows
against the pooled macro) and the sample-complexity slopes with an exactly enumerated
bootstrap. The audit of 18 September found no arithmetic disagreement anywhere; every
problem it found was in what a number was said to mean, which is what a second
computation cannot catch and a second reader can.

\section{The cluster runs}
\label{sec:hpc}

\paragraph{Deploy.} \code{hpc/icml\_deploy.sh} ships a \emph{committed revision}, never
the working tree. It refuses a dirty tree (untracked outputs excepted), writes
\code{git archive HEAD} to a tarball, computes its SHA-256, copies it to the login node,
verifies the digest on the other side, extracts it under a directory named by the short
commit and a timestamp, and writes a \code{REVISION} file (commit, archive digest,
deployment time, dirty list) at the package root, which \code{provenance()} reads on
nodes without \code{.git}. It creates the shared virtual environment once, runs an import
check, and writes \code{hpc/.last\_deploy.env} with the host, root, source and output
paths and the tag, which the submission script sources.

\paragraph{Submit.} \code{hpc/icml\_submit.sh <experiment>} maps a name to the driver
command with its \code{--set} overrides (the same strings that appear in the registry
entries, so that the amended designs of E06c, E07b, E07c, E11b, E12b and E14 are
reproducible from the script), and submits one SLURM job per experiment: one node, one
task, the requested CPUs, memory and time; \code{--propagate=NONE} and an unlimited CPU
time limit, because the login shell's limit had once been inherited and had killed the
workers of two jobs; single-threaded BLAS so that the workers do not oversubscribe the
node; \code{PYTHONPATH} pointing at the deployment; the \code{REVISION} file printed into
the log; \code{--workers} set to the job's CPU count for drivers that declare the flag.
Dependencies are wired with \code{afterok}, so that \code{e03-confirm} waits for
\code{e03-screen} and \code{e05} for both E02a and E03's confirmation; \code{all} submits
the whole registered set with those dependencies.

\paragraph{Fetch.} \code{hpc/icml\_fetch.sh <E> \ldots} copies the CSV and JSON files
(cell files and checkpoints stay on the cluster) into \code{outputs/<E>/} and prints the
next step: \code{registry.py --check} and \code{freeze.py <id>}.

\section{The chain, end to end}
\label{sec:chain}

For a reader who wants to verify one number, the path from the PDF back to the data is:
\begingroup\sloppy
\begin{enumerate}
\item find the macro in the body file (every quantitative statement is one);
\item find its definition in \code{paper/sections/<x>-numbers.tex}, whose header names
  the generator;
\item read the generator under \code{tools/} to see which CSV columns it reads and what
  it computes; the registry entry's \code{predeclared\_analysis} says the same in words
  and was written before the data;
\item confirm the CSV is frozen: its path and digest are in
  \code{experiments/registry/<id>.json}, and \code{python tools/freeze.py --check} agrees;
\item open the CSV's \code{params\_<part>.json} for the commit, the clean-tree flag, the
  archive digest, the host and the resolved settings, and check that the settings agree
  with the registry's design and the protocol;
\item run \code{python tools/check\_manuscript.py} to regenerate the fragment and
  \code{python tools/audit\_claims.py} to recompute the headline from the CSV
  independently;
\item if the number is a comparison between arms, confirm the rows were joined on the
  seed, the size and the bundle hash, which is what makes it paired.
\end{enumerate}
\endgroup
None of these steps requires trusting the author; each is a file, a digest or a script
that a reader can run. That is the sense in which the paper's numbers are checkable
rather than reported.

\where{\begingroup\sloppy The paper, Appendix A3 (frozen configuration, estimand, arms
and interfaces, provenance) and every generated table and macro. Documents:
\code{docs/EXPERIMENTAL\_CONTRACT.md}, \code{docs/REPRODUCIBILITY.md},
\code{docs/EXPERIMENT\_REGISTRY.md}, \code{docs/NEGATIVE\_RESULTS.md} (N-22),
\code{docs/FIGURE\_TABLE\_STANDARDS.md}. Code: \code{experiments/frozen\_config.py},
\code{experiments/protocol.py}, \code{experiments/common.py} (\code{provenance},
\code{run\_cells}), \code{experiments/registry/*.json}, \code{src/markov\_score/utils.py}
(\code{rng\_for}), \code{tools/freeze.py}, \code{tools/registry.py},
\code{tools/provenance\_gate.py}, \code{tools/check\_manuscript.py},
\code{tools/audit\_claims.py}, \code{paper/LITERALS\_ALLOWED.txt},
\code{hpc/icml\_deploy.sh}, \code{hpc/icml\_submit.sh}, \code{hpc/icml\_fetch.sh},
\code{experiments/e14\_regime\_baselines.py} (the guarded generator).\par\endgroup}

\paragraph{References.}
\begingroup\sloppy
B.~A.~Nosek, C.~R.~Ebersole, A.~C.~DeHaven and D.~T.~Mellor, 2018, The preregistration revolution, \emph{Proceedings of the National Academy of Sciences} 115, 2600--2606.
R.~D.~Peng, 2011, Reproducible research in computational science, \emph{Science} 334, 1226--1227.
G.~K.~Sandve, A.~Nekrutenko, J.~Taylor and E.~Hovig, 2013, Ten simple rules for reproducible computational research, \emph{PLoS Computational Biology} 9, e1003285.
National Institute of Standards and Technology, 2015, \emph{Secure Hash Standard (SHS)}, FIPS PUB 180-4.
A.~B.~Yoo, M.~A.~Jette and M.~Grondona, 2003, SLURM: Simple Linux Utility for Resource Management, in \emph{Job Scheduling Strategies for Parallel Processing}, Springer, 44--60.
P.~Glasserman, 2004, \emph{Monte Carlo Methods in Financial Engineering}, Springer (common random numbers, chapter 4).\par\endgroup
